\documentclass[12pt]{article}
\usepackage[utf8]{inputenc}
\usepackage[T1]{fontenc}
\usepackage[margin=1in]{geometry}
\usepackage{setspace}
\usepackage{xcolor}
\usepackage[colorlinks=true, urlcolor=blue!60!black, linkcolor=blue!60!black]{hyperref}
\setlength{\parindent}{2em}
\setlength{\parskip}{1.25em}
\renewcommand{\baselinestretch}{1.0}
\begin{document}
\begin{center}
{\LARGE \textit{Meyer v. Nebraska}, 262 U.S. 390 (1923)\par}
\vspace{0.8em}
{\large Case Brief\par}
\end{center}
\vspace{1em}

\section{Memory Jogger}

In a 7--2 decision, the Supreme Court invalidated Nebraska's post-World War I ban on teaching modern foreign languages to children below the eighth grade, holding that Fourteenth Amendment ``liberty'' protects a teacher's right to teach, a pupil's right to acquire knowledge, and parents' fundamental right to direct the education and upbringing of their children (\href{https://www.courtlistener.com/opinion/100233/meyer-v-nebraska/}{CourtListener}; \href{https://supreme.justia.com/cases/federal/us/262/390/}{Justia}).

\section{Detailed Case Facts}

Amid intense anti-German and nativist sentiment following World War I, the Nebraska Legislature enacted ``An act relating to the teaching of foreign languages in the State of Nebraska'' (Laws 1919, c.~249, approved April 9, 1919), commonly known as the Siman Act (\href{https://supreme.justia.com/cases/federal/us/262/390/}{Justia}; \href{https://en.wikipedia.org/wiki/Meyer_v._Nebraska}{Wikipedia}). Section 1 prohibited any person, individually or as a teacher in any private, denominational, parochial, or public school, from teaching any subject in any language other than English (\href{https://supreme.justia.com/cases/federal/us/262/390/}{Justia}). Section 2 permitted languages other than English to be taught as languages only after a pupil had attained and successfully passed the eighth grade, as evidenced by a certificate from the county superintendent (\href{https://supreme.justia.com/cases/federal/us/262/390/}{Justia}). Violations were misdemeanors punishable by a fine of \$25 to \$100 or up to thirty days in county jail (\href{https://supreme.justia.com/cases/federal/us/262/390/}{Justia}). The Nebraska Supreme Court had previously construed the act in \textit{Nebraska District of Evangelical Lutheran Synod v. McKelvie}, 187 N.W. 927 (Neb. 1922), to permit ``ancient or dead languages'' such as Latin, Greek, and Hebrew, while banning modern foreign languages such as German, French, Spanish, and Italian (\href{https://supreme.justia.com/cases/federal/us/262/390/}{Justia}).

On May 25, 1920, Robert T. Meyer, an instructor at Zion Parochial School---a school maintained by the Zion Evangelical Lutheran Congregation in Hamilton County, Nebraska---taught reading in the German language to Raymond Parpart, a ten-year-old boy who had not passed the eighth grade, using a collection of Biblical stories (\href{https://supreme.justia.com/cases/federal/us/262/390/}{Justia}; \href{https://www.jurist.org/commentary/2023/06/meyer-v-nebraska/}{Ross, JURIST}). The Hamilton County Attorney entered the classroom, discovered the German instruction, and charged Meyer with violating the Siman Act (\href{https://en.wikipedia.org/wiki/Meyer_v._Nebraska}{Wikipedia}).

\section{Procedural History}

Meyer was tried and convicted in the District Court for Hamilton County, Nebraska, and fined \$25 (\href{https://supreme.justia.com/cases/federal/us/262/390/}{Justia}; \href{https://en.wikipedia.org/wiki/Meyer_v._Nebraska}{Wikipedia}). He appealed to the Supreme Court of Nebraska, arguing that the Siman Act violated the Fourteenth Amendment's Due Process Clause. In \textit{Meyer v. State}, 107 Neb. 657, 187 N.W. 100 (1922), the Nebraska Supreme Court affirmed his conviction, holding that the statute was a valid exercise of the state's police power designed to ensure ``that the English language should be and become the mother tongue of all children reared in this state'' and to prevent ``baneful effects'' from immigrant children being educated in foreign tongues (\href{https://supreme.justia.com/cases/federal/us/262/390/}{Justia}). Meyer sought review in the U.S. Supreme Court by writ of error. The case was argued on February 23, 1923, and decided on June 4, 1923, alongside the companion cases consolidated in \textit{Bartels v. Iowa}, 262 U.S. 404 (1923) (\href{https://supreme.justia.com/cases/federal/us/262/390/}{Justia}; \href{https://supreme.justia.com/cases/federal/us/262/404/}{Justia, \textit{Bartels}}).

\section{Judicial Votes}

The Court voted \textbf{7--2} to reverse the judgment of the Nebraska Supreme Court (\href{https://supreme.justia.com/cases/federal/us/262/390/}{Justia}; \href{https://www.jurist.org/commentary/2023/06/meyer-v-nebraska/}{Ross, JURIST}):

\begin{itemize}
\item \textbf{Majority Opinion (7 Justices):} Justice James Clark McReynolds delivered the opinion of the Court, joined by Chief Justice William Howard Taft and Associate Justices Joseph McKenna, Willis Van Devanter, Louis D. Brandeis, Pierce Butler, and Edward Terry Sanford (\href{https://en.wikipedia.org/wiki/Meyer_v._Nebraska}{Wikipedia}; \href{https://supreme.justia.com/cases/federal/us/262/390/}{Justia}).
\item \textbf{Concurrences:} None.
\item \textbf{Dissenting Opinion (2 Justices):} Justice Oliver Wendell Holmes, Jr. authored a dissent, joined by Justice George Sutherland, published in the companion case \textit{Bartels v. Iowa}, 262 U.S. 404, 412--413 (1923), and cross-referenced at the end of \textit{Meyer}, 262 U.S. at 403 (\href{https://supreme.justia.com/cases/federal/us/262/390/}{Justia}; \href{https://supreme.justia.com/cases/federal/us/262/404/}{Justia, \textit{Bartels}}).
\end{itemize}

\section{Holding}

\textbf{Judgment:} Reversed and remanded (7--2) (\href{https://supreme.justia.com/cases/federal/us/262/390/}{Justia}).

\textbf{Rule Announced:} A state statute forbidding, under criminal penalty, the teaching of any modern language other than English in any private, denominational, parochial, or public school to any child who has not passed the eighth grade arbitrarily infringes the liberty guaranteed by the Due Process Clause of the Fourteenth Amendment and exceeds the police power of the state (\href{https://supreme.justia.com/cases/federal/us/262/390/}{Justia}; \href{https://www.courtlistener.com/opinion/100233/meyer-v-nebraska/}{CourtListener}). Liberty under the Fourteenth Amendment may not be interfered with, under the guise of protecting the public interest, by legislative action that is arbitrary or lacks a reasonable relation to a purpose within the competency of the state (\href{https://supreme.justia.com/cases/federal/us/262/390/}{Justia}).

\textbf{Holding vs. Dicta:} The narrow holding invalidates Nebraska's Siman Act as applied to a parochial school instructor teaching German reading to a ten-year-old child, resting on both the teacher's occupational liberty to teach and the parents' right to engage him to instruct their children (\href{https://supreme.justia.com/cases/federal/us/262/390/}{Justia}). Much of Justice McReynolds's famous catalog of Fourteenth Amendment liberties---such as the rights ``to marry, establish a home and bring up children, to worship God according to the dictates of his own conscience, and generally to enjoy those privileges long recognized at common law as essential to the orderly pursuit of happiness by free men''---was broad dicta in 1923, though it later became foundational authority for modern substantive due process privacy and family-autonomy cases (\href{https://supreme.justia.com/cases/federal/us/262/390/}{Justia}; \href{https://wustllawreview.org/wp-content/uploads/2025/11/14_Aceves_Final-11.10-1.pdf}{Wash. U. L. Rev.}). Likewise, the Court's remarks preserving the state's power to compel school attendance, require English instruction in all schools, prescribe public-school curricula, or potentially impose restrictions during wartime emergencies were dicta (\href{https://supreme.justia.com/cases/federal/us/262/390/}{Justia}).

\section{Analysis of Opinions}

\subsection{Majority Opinion (Justice McReynolds)}

Writing for seven Justices, Justice McReynolds situated the case within the \textit{Lochner}-era substantive due process framework, citing \textit{Allgeyer v. Louisiana}, \textit{Lochner v. New York}, and \textit{Adkins v. Children's Hospital}, but expanded the definition of ``liberty'' well beyond freedom from bodily restraint and economic contract to encompass acquiring useful knowledge, marrying, establishing a home, bringing up children, and worshipping freely (\href{https://supreme.justia.com/cases/federal/us/262/390/}{Justia}). McReynolds emphasized that teaching is a ``useful and honorable'' calling essential to the public welfare, that parents bear a natural duty and corresponding right of control over their children's education, and that ``[m]ere knowledge of the German language cannot reasonably be regarded as harmful'' (\href{https://supreme.justia.com/cases/federal/us/262/390/}{Justia}).

While acknowledging the legislature's understandable post-war desire ``to foster a homogeneous people with American ideals,'' McReynolds insisted that ``a desirable end cannot be promoted by prohibited means'' and that constitutional protections extend to non-English speakers as well as native English speakers (\href{https://supreme.justia.com/cases/federal/us/262/390/}{Justia}). Invoking Plato's \textit{Ideal Commonwealth} and ancient Sparta---where children were taken from parents and reared by state guardians to produce ideal citizens---McReynolds declared that such state submergence of the individual is fundamentally alien to American constitutional institutions (\href{https://supreme.justia.com/cases/federal/us/262/390/}{Justia}). Because the statute targeted only modern foreign languages while leaving ancient languages and other subjects untouched, the Court dismissed health-based justifications and held that, in a time of ``peace and domestic tranquility,'' no emergency justified such an arbitrary prohibition (\href{https://supreme.justia.com/cases/federal/us/262/390/}{Justia}). At the same time, the Court carefully reserved state authority to compel school attendance, make reasonable regulations for all schools (including requiring instruction in English), and set public-school curricula (\href{https://supreme.justia.com/cases/federal/us/262/390/}{Justia}).

\subsection{Dissenting Opinion (Justice Holmes, Joined by Justice Sutherland)}

In his dissent published in \textit{Bartels v. Iowa}, 262 U.S. at 412--413, Justice Holmes (joined by Justice Sutherland) expressed ``hesitation and unwillingness'' in parting with the majority, but urged judicial deference to legislative judgment (\href{https://supreme.justia.com/cases/federal/us/262/404/}{Justia, \textit{Bartels}}). Holmes agreed that fostering a common national tongue is a lawful and desirable end. Reasoning that early childhood is when language fluency is formed, Holmes argued that in districts where children hear only Polish, French, or German at home, a legislature could reasonably conclude that requiring children to hear and speak only English at school is a necessary means to ensure English proficiency (\href{https://supreme.justia.com/cases/federal/us/262/404/}{Justia, \textit{Bartels}}). Because the question was one ``upon which men reasonably might differ,'' Holmes concluded that the Nebraska and Iowa statutes did not pass ``the bounds of reason'' or become an ``arbitrary fiat,'' and that the Constitution did not forbid the states from trying the experiment (\href{https://supreme.justia.com/cases/federal/us/262/404/}{Justia, \textit{Bartels}}). Tellingly, however, Holmes agreed with the majority in invalidating Ohio's statute in \textit{Bohning v. Ohio} because it singled out the German language specifically (\href{https://supreme.justia.com/cases/federal/us/262/404/}{Justia, \textit{Bartels}}).

\section{Examples: Future Applications}

\begin{enumerate}
\item \textbf{Same-Side (Invalidated): Mandatory Public School Attendance.} A state enacts a law requiring all children between ages eight and sixteen to attend public schools, effectively outlawing private and parochial elementary schools. Under \textit{Meyer}, this law unconstitutionally interferes with parents' Fourteenth Amendment liberty to direct the upbringing and education of their children and with private schools' occupational and property interests (as the Court held two years later in \textit{Pierce v. Society of Sisters}, 268 U.S. 510 (1925)) (\href{https://www.americanbar.org/groups/public_interest/child_law/resources/child_law_practiceonline/child_law_practice/vol-35/february-2016/parental-rights-cases-to-know/}{ABA}).
\item \textbf{Same-Side (Invalidated): Ban on Private Cultural After-School Programs.} A state prohibits private heritage academies from teaching Mandarin Chinese or Spanish literacy to children under age twelve in order to promote cultural assimilation. Because knowledge of a foreign language is not harmful and no emergency exists, such a prohibition is arbitrary and violates both teachers' occupational liberty and parents' educational autonomy under \textit{Meyer} (\href{https://supreme.justia.com/cases/federal/us/262/390/}{Justia}).
\item \textbf{Opposite-Side (Upheld): State Requirement That All Schools Teach English Proficiency.} A state statute requires every public and private elementary school to include daily instruction in English reading and writing, while leaving private schools free to offer foreign-language courses in addition. \textit{Meyer} expressly affirmed that ``the power of the State to compel attendance at some school and to make reasonable regulations for all schools, including a requirement that they shall give instructions in English, is not questioned'' (\href{https://supreme.justia.com/cases/federal/us/262/390/}{Justia}).
\item \textbf{Opposite-Side (Upheld): State-Prescribed Public School Curriculum.} A state legislature mandates that public elementary schools teach civics and state history and declines to fund elementary foreign-language classes in public schools. \textit{Meyer} explicitly stated that ``the State's power to prescribe a curriculum for institutions which it supports'' is not impaired by its holding (\href{https://supreme.justia.com/cases/federal/us/262/390/}{Justia}).
\item \textbf{Fence-Sitter (Close Case): Mandatory Public-School Opt-Out Restrictions on Controversial Curricula or Homeschooling Oversight.} A state requires homeschooling parents to submit annual standardized test scores and follow state-approved science and social studies benchmarks, or refuses to let parents opt their children out of a required public-school health curriculum. While \textit{Meyer} recognizes parents' fundamental liberty to direct their children's education, it also preserves the state's authority to enforce reasonable educational regulations and set public-school curricula, leaving courts to balance parental autonomy against state civic and child-welfare interests (\href{https://supreme.justia.com/cases/federal/us/262/390/}{Justia}; \href{https://ideaexchange.uakron.edu/cgi/viewcontent.cgi?article=1397&context=akronlawreview}{Ross, Akron L. Rev.}).
\end{enumerate}

\section{Critique}

\textbf{Originalist and Textualist Perspectives:} Originalists and textualists frequently criticize \textit{Meyer} as an unmoored product of the discredited \textit{Lochner} era, noting that Justice McReynolds---one of the Court's staunchest economic substantive due process conservatives---relied on ``liberty'' in the Due Process Clause without demonstrating that parental educational control was grounded in the Clause's original procedural meaning (\href{https://www.repository.law.indiana.edu/cgi/viewcontent.cgi?article=2021&context=ilj}{Indiana L.J.}; \href{https://digital.sandiego.edu/cgi/viewcontent.cgi?article=1071&context=jcli}{Dwyer, J. Contemp. Legal Issues}). Critics argue that \textit{Meyer}'s expansive catalog of unenumerated rights paved the way for modern substantive due process decisions (\textit{Griswold}, \textit{Roe}, \textit{Lawrence}, \textit{Obergefell}) that invite judicial policymaking under vague ``tradition'' or ``liberty'' standards (\href{https://www.repository.law.indiana.edu/cgi/viewcontent.cgi?article=2021&context=ilj}{Indiana L.J.}; \href{https://www.law.nyu.edu/sites/default/files/ECM_PRO_062707.pdf}{NYU L. Rev.}). Some libertarian-leaning originalists, however, defend \textit{Meyer}'s protection of both occupational and family liberty, though often arguing it should be grounded in the Privileges or Immunities Clause or the Ninth Amendment rather than Due Process (\href{https://scholarship.law.georgetown.edu/cgi/viewcontent.cgi?article=1041&context=facpub}{Georgetown Law Scholarship}; \href{https://scholarship.law.bu.edu/cgi/viewcontent.cgi?filename=0&article=1372&context=faculty_scholarship&type=additional}{Oxford Handbook of the U.S. Constitution}).

\textbf{Progressive and Doctrinal Perspectives:} Many civil-liberties scholars celebrate \textit{Meyer} as a landmark defense of cultural pluralism, minority rights, and personal autonomy against majoritarian xenophobia---and note how \textit{Griswold v. Connecticut} later rehabilitated \textit{Meyer} and \textit{Pierce} as First Amendment and privacy precedents after \textit{Lochner} was repudiated (\href{https://www.jurist.org/commentary/2023/06/meyer-v-nebraska/}{Ross, JURIST}; \href{https://www.repository.law.indiana.edu/cgi/viewcontent.cgi?article=2021&context=ilj}{Indiana L.J.}). Conversely, prominent progressive family-law scholars---most famously Barbara Bennett Woodhouse---have offered a revisionist critique, arguing that \textit{Meyer} elevated a patriarchal, common-law conception of ``the child as property'' of parents into a constitutional right, sidelining children's independent interests in education, voice, and protection from parental isolation or neglect (\href{https://scholarship.law.cornell.edu/cgi/viewcontent.cgi?article=1037&context=cjlpp}{Woodhouse, Cornell J.L. \& Pub. Pol'y}; \href{https://digital.sandiego.edu/cgi/viewcontent.cgi?article=1071&context=jcli}{Dwyer, J. Contemp. Legal Issues}).

\textbf{Underlying Competing Values:} \textit{Meyer} exposes an enduring constitutional tension among three competing forces: (1) family privacy, parental autonomy, and cultural pluralism; (2) democratic majoritarianism and the state's interest in fostering civic assimilation and an educated citizenry; and (3) the independent developmental interests of children themselves (\href{https://ideaexchange.uakron.edu/cgi/viewcontent.cgi?article=1397&context=akronlawreview}{Ross, Akron L. Rev.}; \href{https://wustllawreview.org/wp-content/uploads/2025/11/14_Aceves_Final-11.10-1.pdf}{Wash. U. L. Rev.}).

\section{Key Quotations}

\begin{quote}
``Without doubt, it denotes not merely freedom from bodily restraint, but also the right of the individual to contract, to engage in any of the common occupations of life, to acquire useful knowledge, to marry, establish a home and bring up children, to worship God according to the dictates of his own conscience, and generally to enjoy those privileges long recognized at common law as essential to the orderly pursuit of happiness by free men.''
\end{quote}
\noindent --- Justice McReynolds, Opinion of the Court, 262 U.S. at 399 (\href{https://www.courtlistener.com/opinion/100233/meyer-v-nebraska/}{CourtListener}; \href{https://supreme.justia.com/cases/federal/us/262/390/}{Justia}).

\begin{quote}
``The established doctrine is that this liberty may not be interfered with, under the guise of protecting the public interest, by legislative action which is arbitrary or without reasonable relation to some purpose within the competency of the State to effect.''
\end{quote}
\noindent --- Justice McReynolds, Opinion of the Court, 262 U.S. at 399--400 (\href{https://www.courtlistener.com/opinion/100233/meyer-v-nebraska/}{CourtListener}; \href{https://supreme.justia.com/cases/federal/us/262/390/}{Justia}).

\begin{quote}
``That the State may do much, go very far, indeed, in order to improve the quality of its citizens, physically, mentally and morally, is clear; but the individual has certain fundamental rights which must be respected. The protection of the Constitution extends to all, to those who speak other languages as well as to those born with English on the tongue. Perhaps it would be highly advantageous if all had ready understanding of our ordinary speech, but this cannot be coerced by methods which conflict with the Constitution --- a desirable end cannot be promoted by prohibited means.''
\end{quote}
\noindent --- Justice McReynolds, Opinion of the Court, 262 U.S. at 401 (\href{https://www.courtlistener.com/opinion/100233/meyer-v-nebraska/}{CourtListener}; \href{https://supreme.justia.com/cases/federal/us/262/390/}{Justia}).

\begin{quote}
``Although such measures have been deliberately approved by men of great genius, their ideas touching the relation between individual and State were wholly different from those upon which our institutions rest, and it hardly will be affirmed that any legislature could impose such restrictions upon the people of a State without doing violence to both letter and spirit of the Constitution.''
\end{quote}
\noindent --- Justice McReynolds, Opinion of the Court, 262 U.S. at 402 (\href{https://www.courtlistener.com/opinion/100233/meyer-v-nebraska/}{CourtListener}; \href{https://supreme.justia.com/cases/federal/us/262/390/}{Justia}).

\begin{quote}
``I think I appreciate the objection to the law, but it appears to me to present a question upon which men reasonably might differ, and therefore I am unable to say that the Constitution of the United States prevents the experiment's being tried.''
\end{quote}
\noindent --- Justice Holmes, Dissenting (joined by Justice Sutherland), published in \textit{Bartels v. Iowa}, 262 U.S. 404, 412 (1923) (\href{https://www.courtlistener.com/opinion/100233/meyer-v-nebraska/}{CourtListener}; \href{https://supreme.justia.com/cases/federal/us/262/404/}{Justia, \textit{Bartels}}).


\end{document}
